\documentclass[11pt,a4paper]{article}
\usepackage[utf8]{inputenc}
\usepackage[T1]{fontenc}
\usepackage[french]{babel}
\usepackage{amsmath,amssymb}
\usepackage{booktabs}
\usepackage{graphicx}
\usepackage{geometry}
\geometry{margin=2.3cm}
\usepackage{xcolor}
\usepackage{hyperref}
\hypersetup{colorlinks=true,linkcolor=blue,urlcolor=blue}
\usepackage{enumitem}
\usepackage{multirow}

\title{\textbf{R-SuperCerv : supervision par rapports pour la segmentation\\
d'hémorragie intracrânienne au scanner cérébral\\[4pt]
\large Synthèse exhaustive --- données, méthodes, résultats}}
\author{}
\date{\today}

\begin{document}
\maketitle

\begin{abstract}
\noindent Ce document synthétise l'ensemble du travail réalisé sur \textbf{R-SuperCerv},
portage de \textbf{R-Super} (MICCAI 2025, segmentation de tumeurs supervisée par rapports
via \emph{Volume Loss} et \emph{Ball Loss}) à la segmentation de l'\textbf{hémorragie
intracrânienne} (ICH) au scanner cérébral, avec extension à l'hémorragie intraventriculaire
(IVH) et à l'\oe dème péri-lésionnel (PHE). L'objectif scientifique est de mesurer la
\emph{valeur} et surtout les \emph{limites} de la supervision par rapports radiologiques
face à une baseline entraînée sur masques. La conclusion principale est un \textbf{résultat
négatif net et rigoureux} pour l'ICH : un modèle entraîné sur seulement 25 masques généralise
déjà mieux que ce que les rapports peuvent apporter ; la supervision-volume est à la fois
\emph{inutile} (le rapport est un moins bon estimateur du volume que le modèle),
\emph{sans plafond} (le volume exact n'améliore le Dice que de $+0{,}003$) et
\emph{déstabilisante} (collapse en l'absence d'ancre spatiale ; le collapse est évité par
une ancre de type Kervadec sur pixels sûrs d'intensité, mais le résultat reste sous le
baseline). La seule classe présentant une marge réelle est l'\textbf{IVH}.
\end{abstract}

\tableofcontents

% ======================================================================
\section{Contexte et objectif}
\textbf{R-Super} convertit l'information textuelle des comptes rendus (volume, diamètres,
localisation d'une lésion) en supervision par-voxel pour un segmenteur, via deux pertes :
\begin{itemize}[nosep]
  \item \textbf{Volume Loss} : la somme des probabilités prédites dans une région doit
  correspondre au volume rapporté ;
  \item \textbf{Ball Loss} : place une boule (du diamètre rapporté) là où le modèle est le
  plus confiant, et transforme cela en pseudo-masque per-voxel.
\end{itemize}
\textbf{R-SuperCerv} porte cette idée à l'ICH : lésion \emph{hyperdense} (sang frais
$\sim$50--90~HU) sur parenchyme ($\sim$30--40~HU), au scanner sans injection. Cette
différence d'imagerie (l'ICH a une signature d'intensité que les tumeurs abdominales
n'ont pas) est centrale pour toute l'analyse.

\medskip\noindent\textbf{Question directrice.} La supervision par rapports apporte-t-elle
quelque chose quand une baseline sur masques couvre déjà le domaine cible~? Et si non,
\emph{pourquoi}~?

% ======================================================================
\section{Données}
Deux cohortes de nature et de provenance différentes.

\subsection{Cohorte MASQUES (entraînement supervisé) --- RSNA}
\label{sec:masks}
359 paires CT + masque, masque multi-classe \{1=ICH, 2=IVH, 3=PHE\}, dérivé d'un jeu de
segmentation de type RSNA (\texttt{data\_laurent/NIFTI}). Prévalence des labels
(Tab.~\ref{tab:prevalence}). C'est le seul jeu avec des masques par-voxel pour les trois
types de lésion.

\begin{table}[h]\centering
\caption{Prévalence des labels dans les 359 masques RSNA.}
\label{tab:prevalence}
\begin{tabular}{lcc}
\toprule
Label & Cas présents & \% \\
\midrule
ICH (1) & 359 & 100\,\% \\
IVH (2) & 190 & 53\,\% \\
PHE (3) & 345 & 96\,\% \\
\bottomrule
\end{tabular}
\end{table}

\subsection{Cohorte RAPPORTS (supervision faible) --- CHUM}
\label{sec:reports}
201 examens du CHUM, \textbf{rapports radiologiques en français} (privés), avec
\texttt{matched\_volumes} (CT) et \texttt{matched\_segmentations}. Chaque examen existe en
double~: \texttt{\_0} (scanner initial / admission) et \texttt{\_1} (contrôle / suivi).

\paragraph{Découverte critique sur les masques \emph{rapports}.} Le dossier
\texttt{matched\_segmentations} contient~:
\begin{itemize}[nosep]
  \item \textbf{65 vrais masques} (tous \texttt{\_0}, \texttt{uint8}, valeurs \{0,1,2,3\}) ;
  \item \textbf{136 fichiers qui sont en réalité des CT} (tous \texttt{\_1}, \texttt{int16},
  HU $-1024$--3071) --- \emph{pas} des masques.
\end{itemize}
Il n'existe donc \textbf{aucune segmentation pour les suivis \texttt{\_1}}. Cette découverte
a invalidé une première série d'analyses de taille (calculées en partie contre ces
« masques-poubelle » via \texttt{a==1} sur un CT) et a été corrigée : toutes les analyses
de taille utilisent désormais les \textbf{65 masques \texttt{\_0} valides}.

\paragraph{Structures anatomiques.} TotalSegmentator (tâche \emph{brain\_structures},
16 structures ; 12 utilisées : lobes frontal/pariétal/temporal/occipital, insula, noyaux
caudé/lenticulaire, thalamus, capsule interne, cervelet, tronc, ventricule) a été appliqué
sur les 560 CT (359 masques + 201 rapports).

\paragraph{Extraction LLM.} Qwen2.5-72B-Instruct-AWQ (vLLM) extrait, par lésion~:
\emph{type} (ICH/IVH/SAH/SDH/PHE), \emph{taille} (diamètres), \emph{structure},
\emph{latéralité}, \emph{compte}. Un bug de parsing majeur (une seule information par ligne
$\rightarrow$ seul le \emph{type} était lu) a été corrigé (parsing par blocs), puis les
réponses re-parsées hors-ligne (sans ré-inférer le LLM).

\subsection{Cohorte MBH (évaluation de spécificité) --- RSNA}
1980 scanners de classification (577 positifs avec sous-types RSNA + $\sim$1400 sains),
prétraités à l'identique (1939 après exclusion de \textbf{41 cas de fuite} = 37 patients
déjà dans notre entraînement). Sert à mesurer la \textbf{spécificité} (le modèle hallucine-t-il
des lésions sur cerveau sain~?), impossible sur les rapports (l'ICH y est toujours présent).

\subsection{Découpages (splits)}
\label{sec:splits}
\begin{table}[h]\centering
\caption{Découpages train/test. Les splits rapports sont \textbf{patient-level}
(pas de fuite \texttt{\_0}/\texttt{\_1} d'un même patient entre train et test).}
\label{tab:splits}
\begin{tabular}{llcc}
\toprule
Cohorte & Découpage & Train & Test \\
\midrule
Masques (RSNA) & complet & 305 & 54 \\
Masques (RSNA) & sous-ensembles emboîtés & \multicolumn{2}{c}{$S_{25}\subset S_{50}\subset S_{100}\subset 305$} \\
\midrule
Rapports (CHUM) & \texttt{reports\_all} (\_0+\_1) & 171 & 30 \\
Rapports (CHUM) & \texttt{reports\_baseline0} (\_0) & 55 & 10 \\
\bottomrule
\end{tabular}
\end{table}

\noindent Population rapports : 146 patients (10 \emph{seulement \_0}, 81 \emph{seulement \_1},
55 \emph{les deux}). Volumes ICH exploitables (taille connue) : \texttt{\_0} 63/63 (100\,\%),
\texttt{\_1} 82/118 (69\,\% ; 36 suivis « stables/résolus » sans nouvelle mesure). Le split
\texttt{reports\_all} donne 126 cas avec volume exploitable, contre 53 pour \texttt{baseline0}.

\subsection{Prétraitement}
Réorientation RAI $\rightarrow$ rééchantillonnage 1~mm (XY BSpline, Z plus-proche-voisin
séparé) $\rightarrow$ fenêtrage HU $[0,100]$ $\rightarrow$ z-score $\rightarrow$ \texttt{.npz}.
Padding symétrique de chaque axe à $\geq 128$ voxels (taille de patch). \emph{Note} : ce
padding explique un décalage de grille (ex. $Z{=}111 \rightarrow 129$) qui a dû être répliqué
lors de l'extension à 15 classes.

% ======================================================================
\section{Pipeline et méthodes}
\subsection{Deux étages (R-Super)}
\begin{itemize}[nosep]
  \item \textbf{Stage 1} : segmenteur MedFormer 3D entraîné \emph{sur masques seuls}
  (12 structures + canal(x) lésion). Crop $128^3$, bf16, 60 époques $\times$ 250 itérations.
  \item \textbf{Stage 2} : \emph{fine-tuning} avec ajout des rapports (Volume/Ball Loss),
  en gardant les masques (\texttt{balance\_supervision} : chaque batch mélange masques + rapports).
  Crop $96^3$.
\end{itemize}

\subsection{Estimation du volume (ABC/2)}
$V = \tfrac{A\cdot B\cdot C}{2}$ (diamètres en cm). Pour 2 diamètres, le 3\textsuperscript{e}
axe $C$ est estimé. \textbf{Choix data-driven} : le vrai axe cranio-caudal mesuré sur les
masques valides vérifie $C_{\text{réel}}\approx A$ (médiane 36~mm vs $A{=}40$, $B{=}21$),
pas $C{=}B$. Passer de \texttt{equal\_B} (défaut) à \textbf{\texttt{equal\_A}} débiaise
l'estimation : ratio médian estimé/réel $0{,}69 \rightarrow 1{,}07$, corrélation $r=0{,}86$.

% ======================================================================
\section{Résultats}

\subsection{Le sang ICH est hyperdense et séparable}
Sur les 65 masques valides, distribution des HU (Tab.~\ref{tab:hu}). Une fenêtre
« sang » $[44,100]$~HU capte 87\,\% des voxels-lésion (sensibilité) pour 6--11\,\% de fuite
parenchyme ; un seuil à 44~HU par cas donne sens.\ 87\,\% / spéc.\ 94\,\%. L'ICH est plus
dense que son voisinage dans \textbf{65/65} cas (écart médian 28~HU).

\begin{table}[h]\centering
\caption{HU médian par type (CT natif, 65 cas valides). ICH et IVH hyperdenses
(confondables) ; PHE hypodense (séparé).}
\label{tab:hu}
\begin{tabular}{lccc}
\toprule
Type & HU médian & IQR & \% dans $[44,100]$ \\
\midrule
Parenchyme       & 34 & [28,39] & 11\,\% \\
ICH (label 1)    & 62 & [50,69] & 86\,\% \\
IVH (label 2)    & 58 & [46,70] & 80\,\% \\
PHE (label 3)    & 27 & [24,30] &  0\,\% \\
\bottomrule
\end{tabular}
\end{table}

\subsection{Corrélations de taille (corrigées sur masques valides)}
Après correction du bug des masques \texttt{\_1} : volume rapport (ABC/2) vs volume réel
$r=0{,}85$ (Spearman 0,90) ; aire de coupe max ($\tfrac{\pi}{4}AB$) vs aire réelle
$r=0{,}91$ ; la formule d'ellipse seule (diamètres mesurés sur le masque) $r=1{,}00$.
Ces valeurs \emph{réfutent} la première conclusion erronée (« signal-taille faible $r=0{,}35$ »),
artefact des masques-poubelle.

\subsection{Localisation : structure vs hémisphère}
\label{sec:loc}
Le \texttt{chosen\_segment\_mask} de R-Super = la structure TotalSeg citée. Or son
\textbf{containment} de la lésion est catastrophique~: \textbf{9\,\%} (médiane) de l'ICH
tombe dans la structure citée. À l'inverse, l'\textbf{hémisphère} rapporté (cavité
intracrânienne $\cap$ côté) contient l'ICH à \textbf{96--100\,\%}, et la cavité entière
à 100\,\% (Tab.~\ref{tab:containment}). Cause de l'échec des structures : \emph{l'hémorragie
corrompt TotalSegmentator} (le sang déplace/remplace le parenchyme là où est la lésion).

\begin{table}[h]\centering
\caption{Containment de l'ICH selon la région \texttt{o} (médiane / moyenne).}
\label{tab:containment}
\begin{tabular}{lcc}
\toprule
Région \texttt{o} & Médiane & Moyenne \\
\midrule
Structure TotalSeg citée & 9\,\% & --- \\
Union parenchyme (12 struct.) & 51\,\% & --- \\
Hémisphère (cavité $\cap$ côté) & 100\,\% & 96\,\% \\
Cavité intracrânienne entière & 100\,\% & 92\,\% \\
\bottomrule
\end{tabular}
\end{table}

\noindent La \textbf{latéralité} est fournie directement par le LLM (colonne
\texttt{lateralization}), exploitable dans \textbf{97\,\%} des cas ; la concordance
côté-ICH est de \textbf{98\,\%} (60/61), le seul « discordant » étant en réalité un cas
ICH+IVH mélangé (latéralité correcte). La cavité est construite comme
\texttt{fill\_holes(union des 16 structures)}, insensible à l'hémorragie grâce au
remplissage des trous.

\subsection{Pseudo-masque par intensité (ablation)}
\label{sec:pseudo}
Pseudo-masque = plus grosse composante connexe (compacité : rayon inscrit $\geq 3$~mm) de
sang $[50,85]$~HU dans le c\oe ur de la région, matchée au volume $V_r$
(Tab.~\ref{tab:pseudo}). Le seuil brut à 44~HU échoue (Dice 0,24 : falx, sinus,
calcifications, liseré crânien captés) ; le passage à 50~HU + filtre de compacité +
suppression du liseré fait bondir à \textbf{0,73}. L'oracle (meilleure composante possible)
plafonne à 0,74 : la \emph{sélection} est bonne, le reste est du seuillage. Précision 0,85,
recall 0,64. \textbf{Fuite IVH $\sim$3\,\%}, fuite PHE 0\,\% (PHE hypodense).

\begin{table}[h]\centering
\caption{Dice (pseudo-masque intensité vs GT-ICH) $\times$ région $\times$ niveau de
correction (médianes). Le \emph{grow} calibré nuit ; l'hémisphère $\approx$ cerveau
(la sélection par composante fait la localisation).}
\label{tab:pseudo}
\begin{tabular}{lccc}
\toprule
Région \texttt{o} & base (cap 85) & +cap100+compacité & +grow $V_r$ \\
\midrule
Structure (dilatée 10\,mm) & 0,75 & 0,74 & 0,71 \\
Hémisphère & 0,69 & 0,73 & 0,68 \\
Cerveau entier & 0,69 & 0,73 & 0,68 \\
\bottomrule
\end{tabular}
\end{table}

\subsection{Reproduction de la Ball Loss (pilotée par le modèle)}
La Ball Loss ne pose pas qu'une boule : elle prend les \emph{top-$V_r$ voxels de la
probabilité du modèle} dans une boule placée là où il est confiant. Reproduite exactement
(\texttt{isolate\_tumor}) à partir du stage-1 (Tab.~\ref{tab:ball})~: le modèle brut est
\textbf{déjà à 0,88} ; la ball sur la structure (9\,\%) tombe à 0,58 ; sur l'hémisphère
elle remonte à 0,68 mais \textbf{plafonne bien sous le modèle} --- car la contrainte de
\emph{forme sphérique} + top-$V_r$ dégrade une bonne prédiction (les hémorragies sont
irrégulières). Autrement dit, même avec une région parfaite, la ball ne peut pas dépasser
le modèle.

\begin{table}[h]\centering
\caption{Dice vs GT-ICH (63 cas CHUM). La ball plafonne sous le modèle brut ;
l'intensité fait un peu mieux mais reste sous le modèle.}
\label{tab:ball}
\begin{tabular}{lc}
\toprule
Approche & Dice médian \\
\midrule
Modèle stage-1 brut & \textbf{0,88} \\
Ball sur structure (9\,\%) & 0,58 \\
Ball sur hémisphère (100\,\%) & 0,68 \\
Ball sur cavité & 0,68 \\
Pseudo-masque intensité & 0,72 \\
\bottomrule
\end{tabular}
\end{table}

\subsection{Sweep 15 classes et généralisation cross-domaine}
\label{sec:gen}
Extension du stage-1 à 15 classes (12 structures + ICH + IVH + PHE) --- les masques RSNA
ayant déjà les 3 labels, il suffisait de les extraire (pipeline rendu reproductible via
\texttt{build\_ich\_dataset.py --lesion\_labels} et \texttt{make\_ich3\_dataset.py}).
Entraînement d'un sweep $X\in\{25,50,100,305\}$ masques, puis \textbf{évaluation par classe}
sur test RSNA (\emph{in-domaine}) vs GT-rapports CHUM (\emph{hors-domaine},
Tab.~\ref{tab:gen}).

\begin{table}[h]\centering
\caption{Dice médian par classe $\times$ nombre de masques $X$ : RSNA (in-domaine) /
CHUM (hors-domaine). \textbf{ICH} : aucun gap. \textbf{IVH} : gap modéré, mais s'améliore
avec les masques (apprenable). \textbf{PHE} : gros gap qui \emph{grandit} avec $X$.}
\label{tab:gen}
\begin{tabular}{lcccc}
\toprule
Classe & $X{=}25$ & $X{=}50$ & $X{=}100$ & $X{=}305$ \\
\midrule
ICH & 0,84 / 0,86 & 0,89 / 0,87 & 0,90 / 0,88 & \textbf{0,90 / 0,88} \\
IVH & 0,55 / 0,53 & 0,68 / 0,57 & 0,73 / 0,66 & \textbf{0,77 / 0,71} \\
PHE & 0,56 / 0,45 & 0,61 / 0,46 & 0,67 / 0,47 & \textbf{0,69 / 0,46} \\
\bottomrule
\end{tabular}
\end{table}

\noindent\textbf{Gaps RSNA$-$CHUM} : ICH $\approx 0$ ; IVH $+0{,}06$ ; PHE $+0{,}11
\rightarrow +0{,}23$ (croissant). L'ICH \emph{généralise parfaitement} du RSNA au CHUM
malgré le hors-domaine (le sang hyperdense a la même signature partout).

\paragraph{Le gap PHE est un artefact d'annotation.} Le ratio de volume PHE/ICH diffère
massivement entre cohortes : \textbf{RSNA 1,19} (\oe dème annoté large) vs \textbf{CHUM 0,31}
(liseré serré), $p=8{,}5\times10^{-18}$, HU identiques (25 vs 27). Le modèle apprend
« gros \oe dème » sur RSNA et sur-prédit sur CHUM. Ce n'est \emph{pas} une marge exploitable
par les rapports (qui ne quantifient de toute façon pas l'extension de l'\oe dème).

\paragraph{Couverture des rapports.} Parmi les cas \emph{masqués}, le rapport mentionne un
terme intraventriculaire dans \textbf{97\,\%} des cas IVH (localisé : quel ventricule,
hydrocéphalie), et l'\oe dème / effet de masse dans 83\,\% des cas PHE (mais en \emph{effet
de masse} --- mm de déviation --- pas en extension). Détection LLM vs masque : IVH
sens.\ 83\,\% / spéc.\ 100\,\% ; PHE sens.\ 49\,\% / spéc.\ 83\,\% (les radiologues
sous-mentionnent l'\oe dème).

\subsection{Évaluation MBH (spécificité)}
Le stage-1 atteint une spécificité $\sim$95\,\% (peu de faux positifs sur cerveau sain).
Le stage-2 (supervision-rapports via ball) \textbf{sur-segmente} : la sensibilité monte mais
la spécificité s'effondre (ex.\ $X{=}25$ stage2-all : spéc.\ 18,5\,\%). Cohérent avec la
localisation à 9\,\% : la loss remplit le bon volume au mauvais endroit.

\subsection{Expérience « Size-Loss only » (Kervadec, sans $H$)}
\label{sec:size}
Implémentation d'une contrainte de taille simple (Kervadec et al.\ 2019, Eq.~2), \emph{sans}
le terme d'attache $H(S)$, bande $=$ tolérance $\tau$ autour de $V_r$ :
\[
C(V_S)=\begin{cases}(V_S-a)^2 & V_S<a\\(V_S-b)^2 & V_S>b\\0 & \text{sinon}\end{cases}
\qquad a=(1-\tau)V_r,\; b=(1+\tau)V_r .
\]
Normalisation par $(V_S+V_r)^2$ pour borner $[0,1)$ (÷$V_r^2$ seul laissait exploser la
sur-prédiction). Tolérance \textbf{data-driven $\tau=0{,}4$} (couvre 73\,\% des cas ; $\tau=0{,}2$
n'en couvrirait que 44\,\%). Appliquée \emph{sans région} (V\_S sur tout le crop), BCE-fond
désactivée sur le canal ICH-rapport, tête finale seulement. Sweep $X{=}25 \times
\{$\_0, \_0+\_1$\} \times \lambda\in\{0{,}1;0{,}5;1{,}0\}$ = 6 runs.

\paragraph{Résultat : collapse.} 4/6 configs \textbf{effondrent l'ICH à 0,000}
(prédiction éteinte : max $\sigma=0{,}21$), et IVH/PHE (non supervisés) aussi ;
2 survivent dégradés (base0-$\lambda$0,1 : 0,75 ; all-$\lambda$1,0 : 0,60). Baseline X25 :
0,86. \textbf{Mécanisme} : (i) la size-loss vaut 0 dans la bande, où le modèle démarre
$\Rightarrow$ \emph{zéro gradient} sur le canal-lésion des items-rapport ; (ii) BCE-fond
désactivée $\Rightarrow$ aucune ancre spatiale ; (iii) le \emph{weight-decay} rétrécit alors
les poids à chaque pas ; une fois $V_S<a$, la remontée échoue car le signal est faible,
\emph{non localisé} (« augmente le volume quelque part, pas où ») et sans ancre. Cliquet à
sens unique. \textbf{Le terme $H$ de Kervadec (BCE sur quelques pixels) EST l'ancre
manquante.}

\subsection{Étude des volumes et test oracle}
\label{sec:vol}
Pourquoi la size-loss ne peut rien apporter à l'ICH~? Parce que le modèle prédit \emph{déjà}
le volume mieux que le rapport (Tab.~\ref{tab:vol}), et parce que le volume n'est pas le
goulot : seuiller la prédiction au \textbf{volume exact de la GT} n'améliore le Dice que de
$\mathbf{+0{,}003}$ (0,862 $\rightarrow$ 0,868). Le Dice dépend de l'\emph{overlap spatial},
pas du volume total.

\begin{table}[h]\centering
\caption{Étude volumes (63 cas). Le modèle X25 bat le rapport comme estimateur du volume GT.}
\label{tab:vol}
\begin{tabular}{lcc}
\toprule
Estimateur du volume & $r$ avec GT & Erreur abs.\ médiane \\
\midrule
Modèle X25 (seuil 0,5) & \textbf{0,955} & \textbf{1,47 mL} \\
Modèle X25 (soft $\Sigma\sigma$) & 0,957 & --- \\
Rapport (equal\_A) & 0,874 & 4,44 mL \\
\bottomrule
\end{tabular}
\end{table}

\subsection{Size-loss + ANCRE spatiale (le terme $H$ de Kervadec)}
\label{sec:anchor}
Le collapse de la size-loss vient de l'absence d'ancre spatiale. On réintroduit donc le
terme $H$ (BCE partielle sur des pixels sûrs), avec les pixels sûrs issus de l'\textbf{intensité}
(pas d'annotation manuelle). Le pseudo-masque complet $P$ (Sec.~\ref{sec:pseudo}, Dice 0.73)
étant imparfait, la BCE n'est calculée que sur les \textbf{voxels sûrs}, dérivés de $P$ dans
la loss (noyau de dilatation $\approx$ 3~mm) :
\begin{itemize}[nosep]
  \item \textbf{sûr-positif} (c\oe ur) $=$ $P$ \emph{érodé} de $\sim$3~mm : BCE $\to 1$ ;
  \item \textbf{sûr-négatif} (loin) $=$ au-delà de $P$ \emph{dilaté} de $\sim$3~mm : BCE $\to 0$ ;
  \item \textbf{incertain} (anneau-frontière $\sim$6~mm) : \emph{pas de BCE}.
\end{itemize}
Soit $\texttt{known} = 1 - \texttt{incertain}$ et $\texttt{loss\_seg} = \mathrm{BCE}(\sigma, P)
\times \texttt{known}$, en plus de la size-loss $C(V_S)$. Sweep identique
($X{=}25 \times \{$\_0, \_0+\_1$\} \times \lambda$).

\paragraph{Résultat (Tab.~\ref{tab:anchor}, deux temps).}
(1) \textbf{L'ancre marche} : elle \emph{empêche le collapse} (5/6 configs stables contre 2/6
sans ancre ; médiane des survivants 0,76 vs 0,67 en CHUM) --- confirmant que l'absence
d'ancre spatiale était bien la cause du collapse. (2) \textbf{Mais elle n'améliore pas} :
le meilleur config ancré atteint \textbf{0,80 (CHUM)}, toujours \emph{sous} le baseline 0,86,
exactement comme prédit par l'étude volumes (le pseudo-masque à 0,73 est un moins bon
enseignant que le modèle). L'ancre stabilise la descente, elle ne l'inverse pas
(fig.\ \texttt{previews/sizeloss\_anchor\_recap.png}).

\begin{table}[h]\centering
\caption{ICH Dice médian (CHUM) : sans ancre (collapse) vs avec ancre 'zones sûres'.
Baseline X25 = 0,86. L'ancre stabilise (5/6 vs 2/6) mais reste sous le baseline.}
\label{tab:anchor}
\begin{tabular}{lcc}
\toprule
Config ($\lambda$) & sans ancre & avec ancre \\
\midrule
\_0 \ $\lambda$0,1 & 0,79 & 0,75 \\
\_0 \ $\lambda$0,5 & \textbf{0,00} & \textbf{0,80} \\
\_0 \ $\lambda$1,0 & 0,00 & 0,57 \\
\_0+\_1 \ $\lambda$0,1 & 0,00 & 0,76 \\
\_0+\_1 \ $\lambda$0,5 & 0,00 & 0,00 \\
\_0+\_1 \ $\lambda$1,0 & 0,56 & 0,77 \\
\midrule
Non-collapsés (/6) & 2 & \textbf{5} \\
\bottomrule
\end{tabular}
\end{table}

% ======================================================================
\section{Conclusions}
\subsection{ICH : chapitre clos (résultat négatif solide)}
La supervision-volume par rapports est \textbf{morte sur trois fronts indépendants} :
\begin{enumerate}[nosep]
  \item \textbf{Le rapport ne bat pas le modèle} sur le volume ($r$ 0,874 vs 0,955 ;
  erreur 4,44 vs 1,47 mL) --- même avec seulement 25 masques.
  \item \textbf{Le volume, même exact, n'est pas le goulot} ($+0{,}003$ Dice) --- le Dice
  est spatial.
  \item \textbf{La contrainte-volume déstabilise} sans ancre spatiale (collapse) --- le terme
  $H$ de Kervadec est nécessaire ; et \emph{même avec} l'ancre correcte (H sur pixels sûrs
  par intensité), le meilleur résultat plafonne à 0,80 $<$ 0,86 (Sec.~\ref{sec:anchor}).
\end{enumerate}
De plus, la Ball Loss plafonne sous le modèle (contrainte de forme sphérique), et le modèle
\textbf{généralise déjà parfaitement} du RSNA au CHUM sur l'ICH. \textbf{Toutes les méthodes
testées} (ball, volume-only, size-loss$+$ancre) restent \emph{sous} le baseline : on a
épuisé les variantes, avec les bons réglages, et montré \emph{pourquoi} aucune ne dépasse un
modèle entraîné sur 25 masques.

\subsection{Principe général}
\emph{Une supervision faible n'aide que si son signal bat le modèle sur la quantité
supervisée.} Pour une cible « facile » et transférable (sang hyperdense), un modèle
mask-supervisé atteint le plafond avec très peu de masques $\Rightarrow$ rien à gagner.

\subsection{Pistes par classe}
\begin{itemize}[nosep]
  \item \textbf{ICH} : aucune marge. Message scientifique : \emph{quand et pourquoi la
  supervision-rapports n'aide pas}.
  \item \textbf{IVH} : \emph{seule} marge réelle (gap $+0{,}06$, apprenable ; signal-rapport
  fort : présence 97\,\%, localisation ventriculaire ; hyperdense). Prochaine étape :
  refaire la grille diagnostique (le pseudo-label IVH bat-il le modèle sur CHUM~?) avant
  toute implémentation. Ancre spatiale nécessaire = seeds par intensité ($H$) + contrainte.
  \item \textbf{PHE} : gap = artefact d'annotation (RSNA 4$\times$ CHUM) ; les rapports ne
  quantifient pas l'\oe dème. Non exploitable par les rapports.
\end{itemize}

% ======================================================================
\section{Reproductibilité}
Tout le code vit dans le dépôt (rien d'utile dans \texttt{/tmp} ou les données) :
\texttt{dataset\_conversion/build\_ich\_dataset.py} (\texttt{--lesion\_labels}),
\texttt{dataset\_conversion/make\_ich3\_dataset.py} (extension 13$\rightarrow$15 classes),
\texttt{training/size\_loss.py} (contrainte de taille ; intégration \texttt{--loss size\_last}
et \texttt{size\_anchor\_last} dans \texttt{calculate\_loss}),
\texttt{dataset\_conversion/make\_ich\_pseudomasks.py} (pseudo-masques ICH par intensité,
ancre),
\texttt{report\_extraction/report\_to\_rsuper\_metadata.py} (\texttt{--third\_axis equal\_A}),
\texttt{run\_sweep15.sh}, \texttt{run\_sizeloss.sh}, \texttt{run\_sizeloss\_anchor.sh},
\texttt{eval\_ich3\_gen.py}, \texttt{eval\_sizeloss.py}. Les chemins sont paramétrés
(\texttt{RSUPER\_DATA}) pour être rejouables sur d'autres données.

\end{document}
